DGFiP, IRCOM revenus 2024 ; OpenFisca-France ; calculs des auteurs.\par Note : $\bar G$(z) est le poids social moyen, rapporté à la valeur d'un euro public, que la société doit accorder aux foyers dont l'assiette dépasse z pour que le barème soit optimal (Saez, 2001, sans effet revenu). Un $\bar G$ négatif signale un taux au-delà du maximum de recette. $\alpha$(z) = z h(z)/(1 $-$ H(z)) est calculé sur la distribution reconstituée de l'assiette brute jusqu'à 150 000 \euro{}, et fixé à 2,1 au-delà (Cabannes, Houdré et Landais, 2014) ; seule la part de l'assiette qui vient du travail réagit au taux.
